\label{sec:dia-seed}

\subsection{The DIA Seed Formula}

The model federal statute (\textit{Democratic Integrity Act}, \texttt{docs/legal/model-federal-statute.md})
specifies a \textit{statutory seed formula} for algorithmic redistricting. The seed
is the initial partition from which \textsc{bisect}'s recursive bisection begins.
Because different seeds can produce different optimal partitions---the solution space
analysis in B.7 shows that convergence search at $T=600$ produces seat-count variance
$< 0.3$ across 50 seeds, but the specific partition varies---the seed formula is
a critical parameter that must be fixed before redistricting begins.

The DIA seed formula specifies:

\begin{equation}
\text{seed} = \text{SHA-256}(\text{state\_fips} \| \text{year} \| \text{census\_hash}) \mod 2^{64}
\end{equation}

where \texttt{census\_hash} is the SHA-256 hash of the P.L.\ 94-171 redistricting
data file as released by the Census Bureau, \texttt{state\_fips} is the two-digit
FIPS code for the state, and \texttt{year} is the census year. This formula ensures that:

\begin{enumerate}
\item The seed is deterministic: the same census data always produces the same seed.
\item The seed is unpredictable before census data release: no party can compute
  the seed in advance and pre-calculate optimal gerrymandering strategies.
\item The seed is verifiable after census data release: any party can verify the
  computation independently.
\end{enumerate}

\subsection{Why Last-Minute Changes Are Prohibited}

The DIA seed formula makes the seed a function of the census data itself, which is not
available until census data release. This means that the algorithm cannot be fully
configured until the census data arrives. However, everything \textit{except} the seed
can and must be pre-registered: the structure, weights, search strategy, balance tolerance,
and all other configuration parameters.

The legal significance is this: if any configuration parameter other than the seed
is changed after census data release, the redistricting authority must explain why
the change was made. Any change that increases the partisan advantage of the resulting
map triggers a rebuttable presumption of gerrymandering under the DIA's transparency
provisions. Courts have interpreted similar provisions in state redistricting statutes
as creating a burden-shifting framework: changes to algorithmic parameters after
census data is available are presumptively partisan unless the redistricting authority
can demonstrate a non-partisan technical justification.

The practical consequence: if the \textsc{bisect} infrastructure for 2030 is not
pre-registered by July 2030, the redistricting authority will either use un-registered
configurations (risking legal challenge) or delay redistricting past constitutional
deadlines (also risking legal challenge). There is no third option.

\subsection{Configuration Stability and Pre-Registration Strategy}

The B-series algorithm papers demonstrate that the key configuration choices---structure
(\texttt{prime-factor} for most states), weights (\texttt{geographic} for geographic
compactness, \texttt{county} for subdivision preservation), and search (\texttt{convergence}
at $T=600$)---produce stable results across the 2000--2020 redistricting cycles. This
stability is the empirical foundation for pre-registration: we can pre-register
configurations based on 2020 performance and have reasonable confidence that the
same configurations will perform appropriately on 2030 census data.

States with unusual geographic configurations (Montana, Wyoming, Alaska, Hawaii)
require case-by-case configuration review. The Q.2 paper provides a state-by-state
configuration audit for the 2030 pre-registration.

\begin{table}[ht]
\centering
\caption{Pre-Registration Schedule: Key Dates}
\label{tab:registration}
\begin{tabular}{ll}
\toprule
Date & Milestone \\
\midrule
January 2029 & Reapportionment projections finalized; $k$-values locked \\
July 2029 & ACS 2026--2028 data available; demographic validation complete \\
January 2030 & Block-group adjacency construction begins \\
July 2030 & Pre-registration deadline for 2031 cycle (DIA, 6-month rule) \\
December 2030 & LODES 2028 data available; incorporated in registered config \\
January 2031 & Block-group adjacency validation complete; final infrastructure check \\
April 2031 & Census data release; seed computation; redistricting begins \\
\bottomrule
\end{tabular}
\end{table}
